\section{大陸地殻の組成・上下構造}

\subsection{令和2年度 問28}
\label{q:r02-reg-applied-28}

\noindent\textbf{原典：} sources/R02/17\_applied\_science.pdf、PDF物理10頁。

\textbf{公式正答：} 3（single）。

\par\textbf{検討未確定：} 年代・厚さ・密度・面積は代表値と傾向である。二層の均一性や局所的な密度序列、約40\%の厳密な境界定義まで原典は指定しない。

\paragraph{問題の要旨} 大陸地殻の上下の組成と密度、典型厚さ、年代、面積を整理する。

\Needspace{4\baselineskip}\paragraph{確認した条件・記述（独自要約）}\mbox{}\par

\noindent\textbf{確認範囲：} 問題幹、数値、単位、否定、全5候補を原典画像で照合した。\par

\noindent\textbf{成立条件：} 地殻全体の表面積を分母とする。平均と全地点、現存岩石年代と地球誕生、岩石名と平均化学組成を区別する。\par

\paragraph{問われる核心} 大陸地殻の上下の組成と密度、典型厚さ、年代、面積を整理する。

\paragraph{必要知識} 地殻全体の表面積を分母とする。平均と全地点、現存岩石年代と地球誕生、岩石名と平均化学組成を区別する。

\paragraph{考え方・正答を選ぶ理由（査読待ち初稿）} 大陸地殻の上部は平均的に珪長質に富み、下部はそれより苦鉄質に富むというモデルである。下部が主に花崗岩質という3が不適切である。大陸地殻は陸地だけでなく海面下の大陸棚にも分布するため、面積約40\%と陸地約30\%を混同しない。

\Needspace{4\baselineskip}\paragraph{選択肢の検討（査読待ち初稿）}\mbox{}\par

\noindent\textbf{1：条件付きで適切} 岩石の年代は約40億年前から現在。大陸には約40億年前の古い岩石が残り、現在も岩石が形成される。地球の形成年代や最古ジルコンの年代そのものとは区別する。\par

\noindent\textbf{2：条件付きで適切} 平均厚さ30〜40 km。代表的平均値であり、厚い山脈下や薄いリフトを含む全地点の範囲ではない。\par

\noindent\textbf{3：不適切} 下部は主に花崗岩質。一般的な組成モデルで下部は上部より苦鉄質に富む。上部の珪長質組成を下部へ置き換えた記述である。\par

\noindent\textbf{4：条件付きで適切} 密度は上部より下部が大きい傾向。下部の組成・圧力に対応する代表的傾向である。個々の局所岩体を密度だけで機械的に区分しない。\par

\noindent\textbf{5：条件付きで適切} 地殻表面積の約40\%が大陸地殻。大陸棚など海面下の大陸地殻を含む。陸地の面積比約3割とは分母対象が異なる。\par

\paragraph{再利用する考え方} 岩石名・上下・密度を関連付けて読み、40億年、40 km、40\%という似た数字の次元と分母を分ける。

\paragraph{共通理論} \hyperref[theory:interior:continental]{第\ref*{theory:interior:continental}節}, \hyperref[theory:interior:ages]{第\ref*{theory:interior:ages}節}

\par 参考：\cite{umich_crust_teaching,tokyo_crust_gallery}。

\paragraph{関連問題} \hyperref[q:r01-ret-applied-34]{R01-RET-APPLIED-34}, \hyperref[q:r04-reg-applied-26]{R04-REG-APPLIED-26}

\par 難易度 2（校準初稿）、優先度未評価。

\subsection{令和4年度 問26}
\label{q:r04-reg-applied-26}

\noindent\textbf{原典：} sources/R04/17\_applied\_science.pdf、PDF物理9頁。

\textbf{公式正答：} 1（single）。

\par\textbf{検討未確定：} 年代・厚さ・密度・面積は代表値と傾向である。二層の均一性や局所的な密度序列、約40\%の厳密な境界定義まで原典は指定しない。

\paragraph{問題の要旨} 大陸地殻の上下の組成と密度、典型厚さ、年代、面積を整理する。

\Needspace{4\baselineskip}\paragraph{確認した条件・記述（独自要約）}\mbox{}\par

\noindent\textbf{確認範囲：} 問題幹、数値、単位、否定、全5候補を原典画像で照合した。\par

\noindent\textbf{成立条件：} 地殻全体の表面積を分母とする。平均と全地点、現存岩石年代と地球誕生、岩石名と平均化学組成を区別する。\par

\paragraph{問われる核心} 大陸地殻の上下の組成と密度、典型厚さ、年代、面積を整理する。

\paragraph{必要知識} 地殻全体の表面積を分母とする。平均と全地点、現存岩石年代と地球誕生、岩石名と平均化学組成を区別する。

\paragraph{考え方・正答を選ぶ理由（査読待ち初稿）} 大陸地殻の上部は平均的に珪長質に富み、下部はそれより苦鉄質に富むというモデルである。下部が主に花崗岩質という1が不適切である。大陸地殻は陸地だけでなく海面下の大陸棚にも分布するため、面積約40\%と陸地約30\%を混同しない。

\Needspace{4\baselineskip}\paragraph{選択肢の検討（査読待ち初稿）}\mbox{}\par

\noindent\textbf{1：不適切} 下部は主に花崗岩質。一般的な組成モデルで下部は上部より苦鉄質に富む。上部の珪長質組成を下部へ置き換えた記述である。\par

\noindent\textbf{2：条件付きで適切} 岩石の年代は約40億年前から現在。大陸には約40億年前の古い岩石が残り、現在も岩石が形成される。地球の形成年代や最古ジルコンの年代そのものとは区別する。\par

\noindent\textbf{3：条件付きで適切} 平均厚さ30〜40 km。代表的平均値であり、厚い山脈下や薄いリフトを含む全地点の範囲ではない。\par

\noindent\textbf{4：条件付きで適切} 密度は上部より下部が大きい傾向。下部の組成・圧力に対応する代表的傾向である。個々の局所岩体を密度だけで機械的に区分しない。\par

\noindent\textbf{5：条件付きで適切} 地殻表面積の約40\%が大陸地殻。大陸棚など海面下の大陸地殻を含む。陸地の面積比約3割とは分母対象が異なる。\par

\paragraph{再利用する考え方} 岩石名・上下・密度を関連付けて読み、40億年、40 km、40\%という似た数字の次元と分母を分ける。

\paragraph{共通理論} \hyperref[theory:interior:continental]{第\ref*{theory:interior:continental}節}, \hyperref[theory:interior:ages]{第\ref*{theory:interior:ages}節}

\par 参考：\cite{umich_crust_teaching,tokyo_crust_gallery}。

\paragraph{関連問題} \hyperref[q:r02-reg-applied-28]{R02-REG-APPLIED-28}

\par 難易度 2（校準初稿）、優先度未評価。

\section{地殻とモホ}

\subsection{平成26年度 問30}
\label{q:h26-reg-applied-30}

\noindent\textbf{原典：} sources/H26/17\_applied\_science.pdf、PDF物理12頁。

\textbf{公式正答：} 4（single）。

\par\textbf{検討未確定：} 二つの均質な岩石層という記述と30〜40 kmの数値は近似である。原典の一般化を地球全域の厳密な組成・厚さの保証へ拡張しない。

\paragraph{問題の要旨} 地殻の位置・組成・厚さ・生成場所を、マントルとプレートの区分から判定する。

\Needspace{4\baselineskip}\paragraph{確認した条件・記述（独自要約）}\mbox{}\par

\noindent\textbf{確認範囲：} 問題幹、数値、単位、否定、全5候補を原典画像で照合した。\par

\noindent\textbf{成立条件：} 地殻とプレートは別概念である。厚さには海水を含めない。花崗岩層・玄武岩層という二層像は平均的組成の近似として読む。\par

\paragraph{問われる核心} 地殻の位置・組成・厚さ・生成場所を、マントルとプレートの区分から判定する。

\paragraph{必要知識} 地殻とプレートは別概念である。厚さには海水を含めない。花崗岩層・玄武岩層という二層像は平均的組成の近似として読む。

\paragraph{考え方・正答を選ぶ理由（査読待ち初稿）} 海洋地殻は主として発散境界の海嶺でマグマが固まって形成される。海溝は海洋プレートが沈み込む側であり、海溝で形成されるとする4が不適切である。大陸と海洋の組成・厚さの典型的相違はこの形成と再循環の違いに対応する。

\Needspace{4\baselineskip}\paragraph{選択肢の検討（査読待ち初稿）}\mbox{}\par

\noindent\textbf{1：適切} 地殻は大きな層構造の最外層。地殻・マントル・核という区分では最外側に位置する。大気をこの固体地球の区分に含めない。\par

\noindent\textbf{2：条件付きで適切} 大陸は花崗岩層と玄武岩層、海洋は玄武岩層。大陸上部は珪長質寄り、下部はより苦鉄質、海洋地殻は玄武岩・斑れい岩を主体とする。実岩石が全域で二つの均一層を作るという断定ではない。\par

\noindent\textbf{3：条件付きで適切} 大陸30〜40 km、海水を除く深海の地殻は大部分10 km未満。代表的厚さの比較であり、山脈下の厚い地殻や海台などの例外を除く典型像である。\par

\noindent\textbf{4：不適切} 海洋地殻は海溝で形成。新しい海洋地殻の主な生成場所は海嶺である。海溝の付加体形成と海洋地殻の生成を同一視しない。\par

\noindent\textbf{5：適切} 地殻とマントルの境界はモホロビチッチ不連続面。地震波の速度変化で認識される地殻底の境界である。リソスフェア底面とは異なる。\par

\paragraph{再利用する考え方} 岩石組成の区分と力学的な硬さの区分を先に分け、生成する海嶺と沈み込む海溝を対にして覚える。

\paragraph{共通理論} \hyperref[theory:interior:crust]{第\ref*{theory:interior:crust}節}, \hyperref[theory:interior:moho]{第\ref*{theory:interior:moho}節}, \hyperref[theory:interior:boundaries]{第\ref*{theory:interior:boundaries}節}

\par 参考：\cite{usgs_interior_structure,usgs_plate_boundaries,umich_crust_teaching}。

\paragraph{関連問題} \hyperref[q:h28-reg-applied-30]{H28-REG-APPLIED-30}, \hyperref[q:h29-reg-applied-26]{H29-REG-APPLIED-26}, \hyperref[q:r07-reg-applied-24]{R07-REG-APPLIED-24}

\par 難易度 2（校準初稿）、優先度未評価。

\subsection{平成28年度 問30}
\label{q:h28-reg-applied-30}

\noindent\textbf{原典：} sources/H28/17\_applied\_science.pdf、PDF物理12頁。

\textbf{公式正答：} 4（single）。

\par\textbf{検討未確定：} 二つの均質な岩石層という記述と30〜40 kmの数値は近似である。原典の一般化を地球全域の厳密な組成・厚さの保証へ拡張しない。

\paragraph{問題の要旨} 地殻の位置・組成・厚さ・生成場所を、マントルとプレートの区分から判定する。

\Needspace{4\baselineskip}\paragraph{確認した条件・記述（独自要約）}\mbox{}\par

\noindent\textbf{確認範囲：} 問題幹、数値、単位、否定、全5候補を原典画像で照合した。\par

\noindent\textbf{成立条件：} 地殻とプレートは別概念である。厚さには海水を含めない。花崗岩質層・玄武岩質層という二層像は平均的組成の近似として読む。\par

\paragraph{問われる核心} 地殻の位置・組成・厚さ・生成場所を、マントルとプレートの区分から判定する。

\paragraph{必要知識} 地殻とプレートは別概念である。厚さには海水を含めない。花崗岩質層・玄武岩質層という二層像は平均的組成の近似として読む。

\paragraph{考え方・正答を選ぶ理由（査読待ち初稿）} 海洋地殻は主として発散境界の海嶺でマグマが固まって形成される。海溝は海洋プレートが沈み込む側であり、海溝で形成されるとする4が不適切である。大陸と海洋の組成・厚さの典型的相違はこの形成と再循環の違いに対応する。

\Needspace{4\baselineskip}\paragraph{選択肢の検討（査読待ち初稿）}\mbox{}\par

\noindent\textbf{1：適切} 地殻は大きな層構造の最外層。地殻・マントル・核という区分では最外側に位置する。大気をこの固体地球の区分に含めない。\par

\noindent\textbf{2：条件付きで適切} 大陸は花崗岩質層と玄武岩質層、海洋は玄武岩質層。大陸上部は珪長質寄り、下部はより苦鉄質、海洋地殻は玄武岩・斑れい岩を主体とする。実岩石が全域で二つの均一層を作るという断定ではない。\par

\noindent\textbf{3：条件付きで適切} 大陸30〜40 km、海水を除く深海の地殻は大部分10 km未満。代表的厚さの比較であり、山脈下の厚い地殻や海台などの例外を除く典型像である。\par

\noindent\textbf{4：不適切} 海洋地殻は海溝で形成。新しい海洋地殻の主な生成場所は海嶺である。海溝の付加体形成と海洋地殻の生成を同一視しない。\par

\noindent\textbf{5：適切} 地殻とマントルの境界はモホロビチッチ不連続面。地震波の速度変化で認識される地殻底の境界である。リソスフェア底面とは異なる。\par

\paragraph{再利用する考え方} 岩石組成の区分と力学的な硬さの区分を先に分け、生成する海嶺と沈み込む海溝を対にして覚える。

\paragraph{共通理論} \hyperref[theory:interior:crust]{第\ref*{theory:interior:crust}節}, \hyperref[theory:interior:moho]{第\ref*{theory:interior:moho}節}, \hyperref[theory:interior:boundaries]{第\ref*{theory:interior:boundaries}節}

\par 参考：\cite{usgs_interior_structure,usgs_plate_boundaries,umich_crust_teaching}。

\paragraph{関連問題} \hyperref[q:h26-reg-applied-30]{H26-REG-APPLIED-30}, \hyperref[q:h29-reg-applied-26]{H29-REG-APPLIED-26}

\par 難易度 2（校準初稿）、優先度未評価。

\subsection{平成29年度 問26}
\label{q:h29-reg-applied-26}

\noindent\textbf{原典：} sources/H29/17\_applied\_science.pdf、PDF物理9頁。

\textbf{公式正答：} 4（single）。

\par\textbf{検討未確定：} 二つの均質な岩石層という記述と30〜40 kmの数値は近似である。原典の一般化を地球全域の厳密な組成・厚さの保証へ拡張しない。

\paragraph{問題の要旨} 地殻の位置・組成・厚さ・生成場所を、マントルとプレートの区分から判定する。

\Needspace{4\baselineskip}\paragraph{確認した条件・記述（独自要約）}\mbox{}\par

\noindent\textbf{確認範囲：} 問題幹、数値、単位、否定、全5候補を原典画像で照合した。\par

\noindent\textbf{成立条件：} 地殻とプレートは別概念である。厚さには海水を含めない。花崗岩層・玄武岩層という二層像は平均的組成の近似として読む。\par

\paragraph{問われる核心} 地殻の位置・組成・厚さ・生成場所を、マントルとプレートの区分から判定する。

\paragraph{必要知識} 地殻とプレートは別概念である。厚さには海水を含めない。花崗岩層・玄武岩層という二層像は平均的組成の近似として読む。

\paragraph{考え方・正答を選ぶ理由（査読待ち初稿）} 海洋地殻は主として発散境界の海嶺でマグマが固まって形成される。海溝は海洋プレートが沈み込む側であり、海溝で形成されるとする4が不適切である。大陸と海洋の組成・厚さの典型的相違はこの形成と再循環の違いに対応する。

\Needspace{4\baselineskip}\paragraph{選択肢の検討（査読待ち初稿）}\mbox{}\par

\noindent\textbf{1：適切} 地殻は大きな層構造の最外層。地殻・マントル・核という区分では最外側に位置する。大気をこの固体地球の区分に含めない。\par

\noindent\textbf{2：条件付きで適切} 大陸は花崗岩層と玄武岩層、海洋は玄武岩層。大陸上部は珪長質寄り、下部はより苦鉄質、海洋地殻は玄武岩・斑れい岩を主体とする。実岩石が全域で二つの均一層を作るという断定ではない。\par

\noindent\textbf{3：条件付きで適切} 大陸30〜40 km、海水を除く深海の地殻は大部分10 km未満。代表的厚さの比較であり、山脈下の厚い地殻や海台などの例外を除く典型像である。\par

\noindent\textbf{4：不適切} 海洋地殻は海溝で形成。新しい海洋地殻の主な生成場所は海嶺である。海溝の付加体形成と海洋地殻の生成を同一視しない。\par

\noindent\textbf{5：適切} 地殻とマントルの境界はモホロビチッチ不連続面。地震波の速度変化で認識される地殻底の境界である。リソスフェア底面とは異なる。\par

\paragraph{再利用する考え方} 岩石組成の区分と力学的な硬さの区分を先に分け、生成する海嶺と沈み込む海溝を対にして覚える。

\paragraph{共通理論} \hyperref[theory:interior:crust]{第\ref*{theory:interior:crust}節}, \hyperref[theory:interior:moho]{第\ref*{theory:interior:moho}節}, \hyperref[theory:interior:boundaries]{第\ref*{theory:interior:boundaries}節}

\par 参考：\cite{usgs_interior_structure,usgs_plate_boundaries,umich_crust_teaching}。

\paragraph{関連問題} \hyperref[q:h26-reg-applied-30]{H26-REG-APPLIED-30}, \hyperref[q:h28-reg-applied-30]{H28-REG-APPLIED-30}

\par 難易度 2（校準初稿）、優先度未評価。

\section{島弧配置と熱流量}

\subsection{令和元年度再試験 問33}
\label{q:r01-ret-applied-33}

\noindent\textbf{原典：} sources/R01\_retest/17\_applied\_science.pdf、PDF物理11頁。

\textbf{公式正答：} 4（single）。

\par\textbf{検討未確定：} 5島弧の区分と縁海の同義扱いは慣用的である。熱流量の広域傾向と地点ごとの地下水・火山の熱輸送を混同しない。

\paragraph{問題の要旨} 島弧の海溝・火山帯・背弧海盆・熱流量・沈み込む地震面の配置を判定する。

\Needspace{4\baselineskip}\paragraph{確認した条件・記述（独自要約）}\mbox{}\par

\noindent\textbf{確認範囲：} 問題幹、数値、単位、否定、全5候補を原典画像で照合した。\par

\noindent\textbf{成立条件：} 海溝側を前弧、反対の大陸側を背弧とする。日本の5島弧は通常の区分による。地殻熱流量は地表を通る熱流束であり、地下温度そのものではない。\par

\paragraph{問われる核心} 島弧の海溝・火山帯・背弧海盆・熱流量・沈み込む地震面の配置を判定する。

\paragraph{必要知識} 海溝側を前弧、反対の大陸側を背弧とする。日本の5島弧は通常の区分による。地殻熱流量は地表を通る熱流束であり、地下温度そのものではない。

\paragraph{考え方・正答を選ぶ理由（査読待ち初稿）} 東北日本弧の広域的な地殻熱流量は、太平洋側の前弧が低く、日本海側の火山・背弧側が高いという分布である。太平洋側で高く大陸側で低いとする4は、この方向を逆にしている。局所的な地下水の移流や火山近傍の異常を広域傾向と区別する。

\Needspace{4\baselineskip}\paragraph{選択肢の検討（査読待ち初稿）}\mbox{}\par

\noindent\textbf{1：条件付きで適切} 5つの弧状列島から構成。千島・東北日本・西南日本・伊豆小笠原・琉球という通常の区分で読む。島弧をさらに分割する流儀まで固定しない。\par

\noindent\textbf{2：適切} 東北日本弧の太平洋側に6000 m以上の海溝が平行。日本海溝と東北日本弧の沈み込み側の配置に対応する。\par

\noindent\textbf{3：条件付きで適切} 大陸側に縁海・背弧海盆。日本海などの典型配置である。縁海という地理区分と背弧海盆という形成過程の区分は完全な同義ではない。\par

\noindent\textbf{4：不適切} 中部軸に火山帯、東北日本の熱流量は太平洋側で高く大陸側で低い。火山帯の記述と合わせた複合肢のうち、熱流量の東西の大小が逆である。\par

\noindent\textbf{5：適切} 海溝から大陸側へ深くなる深発地震面。沈み込むスラブ内の地震分布が深部へ傾く和達・ベニオフ帯に対応する。\par

\paragraph{再利用する考え方} 前弧・火山フロント・背弧を位置の関係で把握し、地温勾配、熱流束、地震面の深度を別の観測量として読む。

\paragraph{共通理論} \hyperref[theory:interior:arc]{第\ref*{theory:interior:arc}節}, \hyperref[theory:interior:heat]{第\ref*{theory:interior:heat}節}

\par 参考：\cite{gsj_heat_flow1991,usgs_plate_boundaries}。

\par 難易度 2（校準初稿）、優先度未評価。

\section{リソスフェアとアセノスフェア}

\subsection{令和7年度 問24}
\label{q:r07-reg-applied-24}

\noindent\textbf{原典：} sources/R07/17\_applied\_science.pdf、PDF物理9頁。

\textbf{公式正答：} 4（single）。

\par\textbf{検討未確定：} 中心温度5000〜6000℃の推定根拠の精査は未完了である。主要プレート枚数と長期的硬さの区分の限定を保持する。

\paragraph{問題の要旨} 地殻・マントルの組成境界とリソスフェア・アセノスフェアの力学境界を区別する。

\Needspace{4\baselineskip}\paragraph{確認した条件・記述（独自要約）}\mbox{}\par

\noindent\textbf{確認範囲：} 問題幹、数値、単位、否定、全5候補を原典画像で照合した。\par

\noindent\textbf{成立条件：} 地球中心の温度は実験と内部モデルからの推定である。地殻・マントルという層と力学的硬さの層は境界位置が一致しない。主要プレート枚数は細分方法による。\par

\paragraph{問われる核心} 地殻・マントルの組成境界とリソスフェア・アセノスフェアの力学境界を区別する。

\paragraph{必要知識} 地球中心の温度は実験と内部モデルからの推定である。地殻・マントルという層と力学的硬さの層は境界位置が一致しない。主要プレート枚数は細分方法による。

\paragraph{考え方・正答を選ぶ理由（査読待ち初稿）} モホロビチッチ不連続面は地殻とマントルの境界である。リソスフェアには地殻だけでなく最上部マントルが含まれる。したがってリソスフェアとアセノスフェアの境界をモホとする4が不適切である。軟らかく流れるという性質は長い時間尺度の変形であり、マントル全体が液体という意味ではない。

\Needspace{4\baselineskip}\paragraph{選択肢の検討（査読待ち初稿）}\mbox{}\par

\noindent\textbf{1：条件付きで適切} 温度は深部ほど高く中心約5000〜6000℃。中心温度の代表的推定範囲である。全深度の厳密な単調性や直接測定値とはしない。\par

\noindent\textbf{2：適切} 上部マントルに流れやすいアセノスフェア。温度が高く長期的に変形しやすい領域であり、通常は主に固体である。\par

\noindent\textbf{3：適切} その上の硬いリソスフェア。地殻と最上部マントルを含む、長期にも比較的強い層である。\par

\noindent\textbf{4：不適切} 両者の境界はモホ。モホは地殻底であり、リソスフェア底の力学境界ではない。\par

\noindent\textbf{5：条件付きで適切} 10数〜20枚程度のプレートの実体はリソスフェア。主要プレートという典型区分では妥当である。微小プレートを数えると枚数が増える。\par

\paragraph{再利用する考え方} 物質の種類、地震波の速度、変形の時間尺度による区分を分ける。硬い層を地殻だけとみなさない。

\paragraph{共通理論} \hyperref[theory:interior:lithosphere]{第\ref*{theory:interior:lithosphere}節}, \hyperref[theory:interior:moho]{第\ref*{theory:interior:moho}節}

\par 参考：\cite{usgs_lithosphere2002,kagoshima_plate2019}。

\paragraph{関連問題} \hyperref[q:h26-reg-applied-30]{H26-REG-APPLIED-30}

\par 難易度 2（校準初稿）、優先度未評価。

\section{静岩圧・上載荷重}

\subsection{平成27年度 問32}
\label{q:h27-reg-applied-32}

\noindent\textbf{原典：} sources/H27/17\_applied\_science.pdf、PDF物理12頁。

\textbf{公式正答：} 4（single）。

\paragraph{問題の要旨} 上載岩石の密度と厚さからモホ面の静岩圧を求める。

\Needspace{4\baselineskip}\paragraph{確認した条件・記述（独自要約）}\mbox{}\par

\noindent\textbf{確認範囲：} 問題幹、数値、単位、否定、全5候補を原典画像で照合した。\par

\noindent\textbf{成立条件：} 厚さ35 km、平均密度2.8×10\textasciicircum{}3 kg/m\textsuperscript{3}、重力加速度9.8 m/s\textsuperscript{2}。上載荷重のみを考え、間隙流体の影響はないという原典指定を使用する。\par

\paragraph{問われる核心} 上載岩石の密度と厚さからモホ面の静岩圧を求める。

\paragraph{必要知識} 厚さ35 km、平均密度2.8×10\textasciicircum{}3 kg/m\textsuperscript{3}、重力加速度9.8 m/s\textsuperscript{2}。上載荷重のみを考え、間隙流体の影響はないという原典指定を使用する。

\paragraph{考え方・正答を選ぶ理由（査読待ち初稿）} 水平面の単位面積当たりに載る岩石の質量はρhである。上載荷重だけを考えるので \(p=\rho gh=(2.8\times10^3)(9.8)(35\times10^3)=9.604\times10^8\,\mathrm{Pa}=0.9604\,\mathrm{GPa}\)。最も近い1 GPaの4を選ぶ。

\Needspace{4\baselineskip}\paragraph{選択肢の検討（査読待ち初稿）}\mbox{}\par

\noindent\textbf{1：不適切} 1 MPa。9.604×10\textasciicircum{}8 Paより約3桁小さい。\par

\noindent\textbf{2：不適切} 10 MPa。導出圧力より約2桁小さい。\par

\noindent\textbf{3：不適切} 100 MPa。導出圧力0.9604 GPaより約1桁小さい。\par

\noindent\textbf{4：適切} 1 GPa。9.604×10\textasciicircum{}8 Paを10\textasciicircum{}9 Paで割った0.9604 GPaに最も近い。\par

\noindent\textbf{5：不適切} 10 GPa。導出圧力より約1桁大きい。\par

\paragraph{再利用する考え方} kmをmへ直し、密度×加速度×長さの単位がPaになることを確認する。静岩圧と間隙水圧及び有効応力を別にする。

\paragraph{共通理論} \hyperref[theory:interior:pressure]{第\ref*{theory:interior:pressure}節}

\par 参考：\cite{usgs_lithosphere2002}。

\par 難易度 2（校準初稿）、優先度未評価。

\section{海洋底と大陸の比較}

\subsection{令和元年度再試験 問34}
\label{q:r01-ret-applied-34}

\noindent\textbf{原典：} sources/R01\_retest/17\_applied\_science.pdf、PDF物理11頁。

\textbf{公式正答：} 2（single）。

\par\textbf{検討未確定：} 大洋底の範囲・モホ深度の基準・堆積速度の地域比較が不明である。候補3の存在命題には例外の精査が必要であり、公式2だけを根拠に候補3を断定しない。

\paragraph{問題の要旨} 海洋底と大陸を、地殻厚さ・堆積・組成・年代の異なる尺度で比較する。

\Needspace{4\baselineskip}\paragraph{確認した条件・記述（独自要約）}\mbox{}\par

\noindent\textbf{確認範囲：} 問題幹、数値、単位、否定、全5候補を原典画像で照合した。\par

\noindent\textbf{成立条件：} 典型的な外洋底と大陸の比較であり、海底すべてを海洋地殻と同一視しない。モホの深さには海面又は固体表面という基準が必要である。\par

\paragraph{問われる核心} 海洋底と大陸を、地殻厚さ・堆積・組成・年代の異なる尺度で比較する。

\paragraph{必要知識} 典型的な外洋底と大陸の比較であり、海底すべてを海洋地殻と同一視しない。モホの深さには海面又は固体表面という基準が必要である。

\paragraph{考え方・正答を選ぶ理由（査読待ち初稿）} 外洋の遠洋性海底では陸源物質の供給が少なく、堆積が遅いことが典型比較となり公式は2である。海洋地殻は薄く、海嶺生成と沈み込みにより現存海洋底の年代は通常数億年以内であり30億年ではない。花崗岩質岩石の全称否定には海洋斜長花崗岩や大陸棚などの範囲の問題があり、3の不適切性を無条件には確定しない。

\Needspace{4\baselineskip}\paragraph{選択肢の検討（査読待ち初稿）}\mbox{}\par

\noindent\textbf{1：不適切} 大洋底のモホは大陸より深い。典型的海洋地殻は薄いため固体表面からのモホ深度は浅い。海面基準でも典型的大陸地殻より浅いが、基準と地域が未指定である。\par

\noindent\textbf{2：条件付きで適切} 大洋底の堆積速度は大陸より遅い。陸源物質の供給が小さい遠洋域の典型である。大陸縁辺の高い堆積速度や大陸の無堆積・侵食域は例外である。\par

\noindent\textbf{3：解釈要確認} 大洋底に花崗岩質岩石が存在。典型的海洋地殻の主体が玄武岩・斑れい岩という説明だけでは、少量の珪長質岩や沈水した大陸地殻の存在を否定できない。原典の大洋底の範囲を要確認とする。\par

\noindent\textbf{4：不適切} 海洋地殻は大陸地殻より厚い。典型的な厚さの大小が逆である。\par

\noindent\textbf{5：不適切} 大洋底の年代は0〜30億年。プレートの再循環が続く現存外洋底は代表的に数億年以内であり、大陸の古い岩石の年代を当てはめている。\par

\paragraph{再利用する考え方} 一般的な大小比較と存在命題を分ける。「主体でない」ことから「存在しない」を導かない。

\paragraph{共通理論} \hyperref[theory:interior:crust]{第\ref*{theory:interior:crust}節}, \hyperref[theory:interior:ages]{第\ref*{theory:interior:ages}節}

\par 参考：\cite{usgs_plate_boundaries,umich_crust_teaching}。

\paragraph{関連問題} \hyperref[q:r02-reg-applied-28]{R02-REG-APPLIED-28}

\par 難易度 2（校準初稿）、優先度未評価。

\section{プレート境界の相対運動}

\subsection{令和3年度 問31}
\label{q:r03-reg-applied-31}

\noindent\textbf{原典：} sources/R03/17\_applied\_science.pdf、PDF物理11頁。

\textbf{公式正答：} 3（single）。

\par\textbf{検討未確定：} 原典の海嶺高度と低速拡大時の溶融度は典型像であり、超低速域の地域条件は省略される。トランスフォームと非活動断裂帯は本文で区別する。

\paragraph{問題の要旨} 三種のプレート境界とリフト・背弧の応力場を、生成・消費・横ずれで区別する。

\Needspace{4\baselineskip}\paragraph{確認した条件・記述（独自要約）}\mbox{}\par

\noindent\textbf{確認範囲：} 問題幹、数値、単位、否定、全5候補を原典画像で照合した。\par

\noindent\textbf{成立条件：} 海嶺間の活動的なトランスフォーム区間と、その外側の非活動的な断裂帯を区別する。海嶺の高さと部分溶融は典型モデルとして扱う。\par

\paragraph{問われる核心} 三種のプレート境界とリフト・背弧の応力場を、生成・消費・横ずれで区別する。

\paragraph{必要知識} 海嶺間の活動的なトランスフォーム区間と、その外側の非活動的な断裂帯を区別する。海嶺の高さと部分溶融は典型モデルとして扱う。

\paragraph{考え方・正答を選ぶ理由（査読待ち初稿）} 海嶺区間の間にある活動的なトランスフォーム断層は両プレートの相対的な横ずれを担う。正断層は伸張で上盤が相対的に下がる断層であり、3の説明とは異なる。断裂帯全長が現在も活動するプレート境界という意味でもない。

\Needspace{4\baselineskip}\paragraph{選択肢の検討（査読待ち初稿）}\mbox{}\par

\noindent\textbf{1：条件付きで適切} 高温の新しい海洋プレートが低密度で海嶺は2〜3 km高い。熱膨張と浮力による海嶺の隆起を説明する典型像である。特定の地形差の厳密値ではない。\par

\noindent\textbf{2：条件付きで適切} 遅い拡大では溶融度が小さくマグマをほとんど伴わない場合。特に超低速拡大の一部ではマグマ供給が少なくマントル岩が露出し得る。遅い海嶺の全域が無マグマという意味ではない。\par

\noindent\textbf{3：不適切} 海嶺をつなぐ断裂帯で正断層運動しトランスフォームとなる。活動区間での基本運動は横ずれである。周辺に局所的な正断層があることとは別の分類である。\par

\noindent\textbf{4：適切} 大陸発散境界の正断層で伸張・薄化・地溝形成。伸張する地殻の相対沈降とリフト形成を対応させた記述である。\par

\noindent\textbf{5：条件付きで適切} 沈み込み帯は圧縮だが背弧が伸張し拡大する場合。沈み込みの存在は背弧の全域が常に圧縮であることを保証しない。海溝後退等に伴う背弧拡大がある。\par

\paragraph{再利用する考え方} 境界の分類を断層名の連想で決めず、相対速度の方向を見て発散・収束・横ずれを判定する。

\paragraph{共通理論} \hyperref[theory:interior:boundaries]{第\ref*{theory:interior:boundaries}節}, \hyperref[theory:interior:ridge]{第\ref*{theory:interior:ridge}節}

\par 参考：\cite{usgs_plate_boundaries,umich_crust_teaching,noaa_ultraslow2004}。

\par 難易度 2（校準初稿）、優先度未評価。

\section{相対プレート速度}

\subsection{令和元年度 問30}
\label{q:r01-reg-applied-30}

\noindent\textbf{原典：} sources/R01/17\_applied\_science.pdf、PDF物理10頁。

\textbf{公式正答：} 5（single）。

\par\textbf{検討未確定：} 北アメリカ固定は明記されるが比較地点とオイラー極モデルは不明である。大学の解説は日本列島に対する代表速度であり、原典の参照系の厳密な数値検証は未完了である。

\paragraph{問題の要旨} 固定した基準プレートに対する日本付近の二つの海洋プレートの速度の桁を比較する。

\Needspace{4\baselineskip}\paragraph{確認した条件・記述（独自要約）}\mbox{}\par

\noindent\textbf{確認範囲：} 問題幹、数値、単位、否定、全5候補を原典画像で照合した。\par

\noindent\textbf{成立条件：} 原典は北アメリカプレート固定を明記する。剛体プレートの相対角速度と地点の位置から速度が決まるため、地点と運動モデルの省略を残す。\par

\paragraph{問われる核心} 固定した基準プレートに対する日本付近の二つの海洋プレートの速度の桁を比較する。

\paragraph{必要知識} 原典は北アメリカプレート固定を明記する。剛体プレートの相対角速度と地点の位置から速度が決まるため、地点と運動モデルの省略を残す。

\paragraph{考え方・正答を選ぶ理由（査読待ち初稿）} 日本付近の代表的沈み込み速度は太平洋プレートが年8〜9 cm、フィリピン海プレートが年4〜5 cmの桁である。候補では9 cm/年と4.5 cm/年の5が対応する。北アメリカ固定という原典の基準を保持し、一般的な日本列島に対する値と完全に同じモデルと証明したとは扱わない。

\Needspace{4\baselineskip}\paragraph{選択肢の検討（査読待ち初稿）}\mbox{}\par

\noindent\textbf{1：不適切} 太平洋6.5 mm/年、フィリピン海2 mm/年。両方が代表的cm/年の運動より一桁程度小さい。\par

\noindent\textbf{2：不適切} 太平洋8 mm/年、フィリピン海7 mm/年。mm/年とcm/年の違いにより代表的な沈み込み速度に達しない。\par

\noindent\textbf{3：不適切} 太平洋2.5 cm/年、フィリピン海1.5 cm/年。両方とも日本付近の代表的速度より小さい組である。\par

\noindent\textbf{4：不適切} 太平洋3 cm/年、フィリピン海8.5 cm/年。太平洋の方が大きい代表比較が逆になっている。\par

\noindent\textbf{5：条件付きで適切} 太平洋9 cm/年、フィリピン海4.5 cm/年。代表的な大きさと大小関係に対応する。原典にない地点・モデルを補って精密値としない。\par

\paragraph{再利用する考え方} cm/年は絶対値ではなく参照系と場所を伴う速度である。拡大速度、沈み込み速度、陸上GNSS点の速度を同じものとして比較しない。

\paragraph{共通理論} \hyperref[theory:interior:velocity]{第\ref*{theory:interior:velocity}節}

\par 参考：\cite{kagoshima_plate2019}。

\par 難易度 2（校準初稿）、優先度未評価。
